\documentclass[11pt]{article}
\newcommand{\ListaTitulo}{Lista 07: Teorema de Bayes}
% Preâmbulo compartilhado: MATF14, Listas de Exercícios 2026
% Usado por ListaNN.tex e GabaritoNN.tex via % Preâmbulo compartilhado: MATF14, Listas de Exercícios 2026
% Usado por ListaNN.tex e GabaritoNN.tex via % Preâmbulo compartilhado: MATF14, Listas de Exercícios 2026
% Usado por ListaNN.tex e GabaritoNN.tex via \input{preamble}

\usepackage[utf8]{inputenc}
\usepackage[T1]{fontenc}
\usepackage[brazil]{babel}
\usepackage{fullpage}
\usepackage{amsmath,amssymb}
\usepackage{enumitem}
\usepackage{xcolor}
\usepackage{fancyhdr}
\usepackage{titlesec}
\usepackage{listings}
\usepackage{hyperref}
\usepackage{booktabs}
\usepackage{graphicx}
\usepackage{tikz}

\definecolor{matf14azul}{HTML}{0D3B54}
\definecolor{matf14dourado}{HTML}{C9971E}

\titleformat{\section}{\normalfont\Large\bfseries\color{matf14azul}}{\thesection}{1em}{}
\titleformat{\subsection}{\normalfont\large\bfseries\color{matf14azul}}{\thesubsection}{1em}{}

\providecommand{\ListaTitulo}{Lista de Exercícios}

\setlength{\headheight}{14pt}
\pagestyle{fancy}
\fancyhf{}
\lhead{\small MATF14: Estatística Econômica I}
\rhead{\small \ListaTitulo}
\cfoot{\thepage}
\renewcommand{\headrulewidth}{0.4pt}

\lstset{
  basicstyle=\ttfamily\small,
  breaklines=true,
  frame=single,
  columns=fullflexible,
  backgroundcolor=\color{gray!8},
  commentstyle=\color{matf14dourado},
  keywordstyle=\color{matf14azul}\bfseries,
  language=R
}

\newcounter{questaocounter}
\newenvironment{questao}[1]{%
  \stepcounter{questaocounter}%
  \bigskip\noindent\textbf{Questão #1.}\ %
}{\par}

\newcommand{\cabecalho}[2]{%
  \begin{center}
    {\Large\bfseries\color{matf14azul} #1}\\[0.3em]
    {\large #2}
  \end{center}
  \vspace{0.5em}
  \noindent\rule{\linewidth}{0.6pt}
  \vspace{0.5em}
}

\newenvironment{solucao}{%
  \par\smallskip\noindent\textit{\color{matf14dourado!80!black}Solução.}\ %
}{\hfill$\blacksquare$\par}


\usepackage[utf8]{inputenc}
\usepackage[T1]{fontenc}
\usepackage[brazil]{babel}
\usepackage{fullpage}
\usepackage{amsmath,amssymb}
\usepackage{enumitem}
\usepackage{xcolor}
\usepackage{fancyhdr}
\usepackage{titlesec}
\usepackage{listings}
\usepackage{hyperref}
\usepackage{booktabs}
\usepackage{graphicx}
\usepackage{tikz}

\definecolor{matf14azul}{HTML}{0D3B54}
\definecolor{matf14dourado}{HTML}{C9971E}

\titleformat{\section}{\normalfont\Large\bfseries\color{matf14azul}}{\thesection}{1em}{}
\titleformat{\subsection}{\normalfont\large\bfseries\color{matf14azul}}{\thesubsection}{1em}{}

\providecommand{\ListaTitulo}{Lista de Exercícios}

\setlength{\headheight}{14pt}
\pagestyle{fancy}
\fancyhf{}
\lhead{\small MATF14: Estatística Econômica I}
\rhead{\small \ListaTitulo}
\cfoot{\thepage}
\renewcommand{\headrulewidth}{0.4pt}

\lstset{
  basicstyle=\ttfamily\small,
  breaklines=true,
  frame=single,
  columns=fullflexible,
  backgroundcolor=\color{gray!8},
  commentstyle=\color{matf14dourado},
  keywordstyle=\color{matf14azul}\bfseries,
  language=R
}

\newcounter{questaocounter}
\newenvironment{questao}[1]{%
  \stepcounter{questaocounter}%
  \bigskip\noindent\textbf{Questão #1.}\ %
}{\par}

\newcommand{\cabecalho}[2]{%
  \begin{center}
    {\Large\bfseries\color{matf14azul} #1}\\[0.3em]
    {\large #2}
  \end{center}
  \vspace{0.5em}
  \noindent\rule{\linewidth}{0.6pt}
  \vspace{0.5em}
}

\newenvironment{solucao}{%
  \par\smallskip\noindent\textit{\color{matf14dourado!80!black}Solução.}\ %
}{\hfill$\blacksquare$\par}


\usepackage[utf8]{inputenc}
\usepackage[T1]{fontenc}
\usepackage[brazil]{babel}
\usepackage{fullpage}
\usepackage{amsmath,amssymb}
\usepackage{enumitem}
\usepackage{xcolor}
\usepackage{fancyhdr}
\usepackage{titlesec}
\usepackage{listings}
\usepackage{hyperref}
\usepackage{booktabs}
\usepackage{graphicx}
\usepackage{tikz}

\definecolor{matf14azul}{HTML}{0D3B54}
\definecolor{matf14dourado}{HTML}{C9971E}

\titleformat{\section}{\normalfont\Large\bfseries\color{matf14azul}}{\thesection}{1em}{}
\titleformat{\subsection}{\normalfont\large\bfseries\color{matf14azul}}{\thesubsection}{1em}{}

\providecommand{\ListaTitulo}{Lista de Exercícios}

\setlength{\headheight}{14pt}
\pagestyle{fancy}
\fancyhf{}
\lhead{\small MATF14: Estatística Econômica I}
\rhead{\small \ListaTitulo}
\cfoot{\thepage}
\renewcommand{\headrulewidth}{0.4pt}

\lstset{
  basicstyle=\ttfamily\small,
  breaklines=true,
  frame=single,
  columns=fullflexible,
  backgroundcolor=\color{gray!8},
  commentstyle=\color{matf14dourado},
  keywordstyle=\color{matf14azul}\bfseries,
  language=R
}

\newcounter{questaocounter}
\newenvironment{questao}[1]{%
  \stepcounter{questaocounter}%
  \bigskip\noindent\textbf{Questão #1.}\ %
}{\par}

\newcommand{\cabecalho}[2]{%
  \begin{center}
    {\Large\bfseries\color{matf14azul} #1}\\[0.3em]
    {\large #2}
  \end{center}
  \vspace{0.5em}
  \noindent\rule{\linewidth}{0.6pt}
  \vspace{0.5em}
}

\newenvironment{solucao}{%
  \par\smallskip\noindent\textit{\color{matf14dourado!80!black}Solução.}\ %
}{\hfill$\blacksquare$\par}


\begin{document}

\cabecalho{Lista de Exercícios 07}{Teorema da probabilidade total e Teorema de Bayes (Cap. 2, Aula 14)}

\noindent \textit{Justifique todas as respostas. As resoluções são discutidas em sala e nos horários de atendimento; não são publicadas neste material.}

\begin{questao}{1} \textbf{(Probabilidade total: inadimplência em carteira segmentada)}

Um banco tem sua carteira de crédito dividida em três segmentos: Pessoa Física (60\% da carteira,
taxa de inadimplência histórica de 4\%), Pequena Empresa (30\% da carteira, taxa de 7\%) e Grande
Empresa (10\% da carteira, taxa de 1\%).

\begin{enumerate}[label=(\alph*)]
  \item Calcule a taxa de inadimplência esperada para a carteira como um todo, usando o teorema da
    probabilidade total.
  \item O gerente de risco observa que a inadimplência geral realizada no trimestre foi de 6\%,
    acima do esperado em (a). Isso necessariamente significa que as taxas por segmento pioraram?
    Cite outra explicação possível envolvendo a composição da carteira.
\end{enumerate}
\end{questao}

\begin{questao}{2} \textbf{(Bayes: origem de um sinistro de seguro)}

Uma seguradora de automóveis tem apólices de dois perfis: "Perfil A" (motoristas com mais de 5
anos de habilitação, 75\% das apólices, probabilidade de sinistro no ano de 5\%) e "Perfil B"
(motoristas novatos, 25\% das apólices, probabilidade de sinistro no ano de 20\%).

\begin{enumerate}[label=(\alph*)]
  \item Calcule a probabilidade de uma apólice sorteada aleatoriamente registrar sinistro no ano.
  \item Dado que uma apólice registrou sinistro, calcule a probabilidade de ela pertencer ao
    Perfil B.
  \item Compare o resultado de (b) com a proporção original de apólices do Perfil B na carteira
    (25\%). O que essa comparação mostra sobre o efeito de observar "houve sinistro" na crença
    sobre o perfil do segurado?
\end{enumerate}
\end{questao}

\begin{questao}{3} \textbf{(Bayes com três categorias: fraude em declarações de imposto de renda)}

A Receita Federal seleciona declarações de imposto de renda para malha fina usando um algoritmo.
Historicamente, 85\% das declarações são "Normais" (probabilidade de sinalização pelo algoritmo:
2\%), 12\% são "Atípicas, mas legítimas" (probabilidade de sinalização: 30\%), e 3\% são
"Fraudulentas" (probabilidade de sinalização: 90\%).

\begin{enumerate}[label=(\alph*)]
  \item Calcule a probabilidade de uma declaração sorteada aleatoriamente ser sinalizada pelo
    algoritmo.
  \item Dado que uma declaração foi sinalizada, calcule a probabilidade de ela ser, respectivamente,
    Normal, Atípica-legítima ou Fraudulenta.
  \item Some as três probabilidades do item (b). O resultado deveria dar 1? Por quê?
\end{enumerate}
\end{questao}

\begin{questao}{4} \textbf{(Discussão: base rate e política pública)}

Um teste de triagem para uma doença rara (prevalência de 0,5\% na população) tem 98\% de
sensibilidade (detecta a doença quando ela existe) e 90\% de especificidade (dá negativo quando a
pessoa é saudável).

\begin{enumerate}[label=(\alph*)]
  \item Calcule $P(\text{doente}\mid\text{teste positivo})$.
  \item Um gestor de saúde pública, ao ver que o teste "acerta 98\% das vezes" (sensibilidade),
    propõe usá-lo para triagem em massa de toda a população, com tratamento imediato para todo
    resultado positivo. Usando o resultado de (a), explique o problema dessa proposta.
\end{enumerate}
\end{questao}

\bigskip
\noindent\textit{A questão a seguir não tem resposta única e não é cobrada em prova no mesmo
formato: fecha a Unidade 2 e é para discussão em grupo ou por escrito, antes da próxima aula.}

\begin{questao}{5 (discussão: síntese da Unidade 2)} \textbf{(Sem resposta única)}

Volte às cinco perguntas de abertura da Aula 01 do semestre (investir em fundo, pesquisa eleitoral,
inflação "dentro da meta", escore de crédito, seguro para condutores jovens). Escolha uma delas e
escreva, em meia página, como você a atacaria usando o vocabulário completo da Unidade 2,
espaço amostral, probabilidade condicional, independência, teorema da probabilidade total e/ou
Teorema de Bayes. Não é preciso resolver numericamente: o objetivo é identificar qual ferramenta se
aplica e justificar a escolha.
\end{questao}


\bigskip
\section*{Exercícios complementares}

\begin{enumerate}[label=\textbf{C\arabic*.}, leftmargin=*]
\item Três máquinas, A, B e C, produzem 60\%, 30\% e 10\% das peças de uma fábrica, com 3\%, 4\% e 5\% de peças defeituosas, respectivamente. Uma peça é sorteada.
\begin{enumerate}[label=(\alph*)]
  \item Qual a probabilidade de ser defeituosa?
  \item Se é defeituosa, qual a probabilidade de ter sido produzida pela máquina C?
\end{enumerate}
\item A urna A tem 8 bolas brancas e 2 pretas; a urna B, 8 pretas e 2 brancas. Uma pessoa escolhe uma urna ao acaso e retira, sem reposição, 2 bolas. Sabendo que as duas são pretas, qual a probabilidade de ter escolhido a urna A?
\item Num país há dois bancos, Alfa e Beta. Alfa tem três vezes mais clientes que Beta; 20\% dos clientes de Alfa e 10\% dos de Beta são aposentados. Um cliente é sorteado.
\begin{enumerate}[label=(\alph*)]
  \item Qual a probabilidade de ser do banco Beta?
  \item Sabendo que é aposentado, qual a probabilidade de ser do banco Beta?
\end{enumerate}
\item Um exame de sangue detecta uma doença em 95\% dos doentes e dá falso positivo em 1\% dos saudáveis. Se 0,5\% da população tem a doença, qual a probabilidade de uma pessoa com resultado positivo estar doente? Refaça supondo que o exame é aplicado só a pessoas com sintomas, das quais 20\% têm a doença.

\end{enumerate}

\bigskip
\needspace{12\baselineskip}
\section*{Aprofundamento: novos tópicos e dados reais}
\noindent\textit{Exercícios sobre os tópicos acrescentados ao Capítulo 2 do livro e às aulas.}

\begin{enumerate}[label=\textbf{A\arabic*.}, leftmargin=*]

\item \textbf{(Probabilidade total com dados reais.)} A força de trabalho do Brasil no 2º trimestre de 2026 (mil pessoas; IBGE, PNAD Contínua):
\begin{center}
\begin{tabular}{l|ccccc}
\toprule
 & Norte & Nordeste & Sudeste & Sul & Centro-Oeste \\
\midrule
Força de trabalho & 8.723 & 25.161 & 48.417 & 17.289 & 9.329 \\
Desocupados & 531 & 1.912 & 2.508 & 552 & 358 \\
\bottomrule
\end{tabular}
\end{center}
\begin{enumerate}[label=(\alph*)]
  \item Calcule a taxa de desocupação de cada região e a participação de cada região na força de trabalho.
  \item Use o teorema da probabilidade total para obter a taxa de desocupação do Brasil.
  \item Uma pessoa desocupada é sorteada. Qual a probabilidade de ela morar no Nordeste? Compare com a participação do Nordeste na força de trabalho.
\end{enumerate}

\item \textbf{(Valor preditivo e prevalência.)} Um teste tem sensibilidade 0,90 e especificidade 0,95. Calcule o valor preditivo positivo
quando a prevalência da condição é 1\%, 10\% e 30\%. O que acontece com a utilidade do teste quando a condição é rara?

\item \textbf{(Bayes na forma de chances.)} A inadimplência de uma carteira é 4\%. Um modelo classifica como ``alto risco'' 85\% dos futuros
inadimplentes e 12\% dos bons pagadores.
\begin{enumerate}[label=(\alph*)]
  \item Calcule a chance a priori de inadimplência e a razão de verossimilhança do sinal ``alto risco''.
  \item Obtenha a chance e a probabilidade a posteriori de inadimplência para um cliente classificado como alto risco.
  \item Um segundo modelo, independente do primeiro dada a situação real e com a mesma precisão, também classifica o cliente como alto risco.
    Qual passa a ser a probabilidade de inadimplência?
\end{enumerate}

\item \textbf{(Transições no mercado de trabalho.)} A cada trimestre, um ocupado perde o emprego com probabilidade $s$ e um desocupado
encontra emprego com probabilidade $f$. Ignorando entradas e saídas da força de trabalho, a taxa de desocupação de longo prazo é
$u^* = s/(s + f)$.
\begin{enumerate}[label=(\alph*)]
  \item Mostre que, com $u^*$, o número de pessoas que entram no desemprego em um trimestre é igual ao número das que saem.
  \item Calcule $u^*$ para $s = 0{,}03$ e $f = 0{,}25$, e depois para $f = 0{,}35$. Interprete.
  \item $s$ e $f$ são probabilidades condicionais. Escreva cada uma na notação $P(\cdot \mid \cdot)$.
\end{enumerate}

\end{enumerate}

\end{document}
